\section{Objectives of the Thesis}

The aim of this thesis is to systematically document, analyze, and compare two contrasting approaches to AI-agent-assisted software development --- one unstructured and one structured --- based on a practical project, and to derive from this comparison concrete, actionable best practices for human-agent collaboration in software engineering.

More specifically, this thesis sets out to:

\begin{enumerate}
    \item \textbf{Document the unstructured approach (Part 1).} Describe how the software was developed using a single-agent workflow with minimal architectural guidance, and analyze the resulting architectural limitations, the influence of prompt precision, and the extent to which human supervision was required.
    \item \textbf{Document the structured approach (Part 2).} Describe the design and implementation of a structured, multi-agent development methodology, including the AI Development Charter, parallel role-specialized agents, mandatory design reviews, the three-tier testing strategy, and the automation of OpenAPI specifications to prevent documentation drift.
    \item \textbf{Evaluate each approach.} Define evaluation criteria and assess each development approach with respect to code quality, test validity, architectural consistency, documentation consistency, and development efficiency.
    \item \textbf{Compare the two approaches.} Conduct a direct cross-approach comparison to determine how the structured methodology affected the limitations observed during the unstructured development phase, including differences in feature implementation, code and test quality, architectural consistency, and documentation reliability.
    \item \textbf{Derive best practices.} Bring the findings together into a set of best practices for human-agent collaboration, covering prompt precision, the role of human oversight, architectural guardrails, recommendations for multi-agent workflows, and strategies for preventing documentation drift.
\end{enumerate}

Taken together, these objectives position the thesis as both an empirical case study of AI-agent-assisted software development and a practically oriented source of recommendations that developers and organizations can draw on when structuring their own collaboration with AI coding agents.

